% !TITLE 野田黄雀行
% !AUTHOR 曹植
% !DYNASTY 魏晋
% !GENRE 五言古诗
% !ORDER 7

\begin{poem}
高树多悲风\textsuperscript{1}，海水扬其波。
利剑不在掌，结友何须多？
不见篱间雀，见鹞\textsuperscript{2}自投罗\textsuperscript{3}？
罗家得雀喜，少年见雀悲。
拔剑捎\textsuperscript{4}罗网，黄雀得飞飞。
飞飞摩\textsuperscript{5}苍天，来下谢少年。
\end{poem}

\begin{annotations}
\notetext{1}{悲风：凄凉的秋风。古人常用“悲风”来渲染凄凉萧瑟的氛围。}
\notetext{2}{鹞：鹞（yào），一种善于捕食小鸟的猛禽。}
\notetext{3}{罗：捕鸟的网。}
\notetext{4}{捎：捎（shāo），这里指拔剑砍断、削去（罗网）。}
\notetext{5}{摩：贴近、迫近。摩苍天，形容黄雀飞得极高，好像贴着青天。}
\end{annotations}

\begin{appreciation}
《野田黄雀行》是曹植后期的诗，像个寓言：黄雀被网住，少年拔剑砍断网，黄雀飞走了。但读过《白马篇》再读这首，就知道这不是寓言，是曹植在说自己。

高树多悲风，海水扬其波——高树上风大，大海里浪高。起兴两句先造气氛：树高了招风，人出了名招祸。利剑不在掌，结友何须多？——剑不在自己手里，交那么多朋友有什么用？这是全诗最沉痛的一句。曹植当时的处境：哥哥曹丕当了皇帝，他的朋友丁仪、丁廙兄弟被杀，他眼睁睁看着，救不了。

不见篱间雀，见鹞自投罗？——你没看见篱笆间的麻雀吗？被鹞鹰一追，慌不择路，自己撞进网里。这是写他朋友的遭遇：不是不机灵，是逃无可逃。罗家得雀喜，少年见雀悲——设网的人捉到雀很高兴，少年看着雀很悲伤。一个喜，一个悲，对比里全是分量。拔剑捎罗网，黄雀得飞飞——少年拔剑削断罗网，黄雀得以飞走。这是全诗最亮的一笔：再坏的世界里，总还有人愿意拔剑。

飞飞摩苍天，来下谢少年——黄雀贴着青天飞了一圈，又飞下来感谢少年。结尾是明亮的，却亮得发冷：现实里曹植的朋友没有得救。诗里写得越圆满，现实越残酷。

还记得《白马篇》吗？同一个曹植，前一篇里少年"捐躯赴国难，视死忽如归"，多么痛快；这一篇里少年只有一把剑，砍的是一张网。中间隔着的，是他哥哥登基。白马篇里的少年是曹植想成为的人；黄雀行里拔剑的少年，是他盼着出现的人——他自己手里没有剑。

交朋友这事，送你一句现成的：不必多。朋友多不多，不看你通讯录里存了多少人，看你真出事的时候，谁还愿意为你拔剑。
\end{appreciation}
